  Implementing and stimulating a project SIEM system allows us the opportunity to understand what the security firm aspect of an actual company goes through. The security team is tasked to handle the internal security framework and network that a company with a security policy heavily depends on when it comes to protecting confidential assets. The biggest assets that needs security protection the most iare employees/customers and so SIEM systems are the most common tools/platforms that helps the security team monitory and regulates events that could show indications of compromise in a system. There are handful of real enterprise level (as well as personal level ) SIEMS, such as Splunk Enterprise/Microsoft Sentinel for commercial level or Graylog/ELK stack/Wazuh for open-source level that are fully published and legally deployed that the avergae citizen can aquire and use. The problem rises when it comes to user management, control, understanding, and overall SIEM access which typically involves some sort of costs; say Splunk Enterprise requires user subscription for SIEM use and deployment.

 There are a handful of higly recommended free-open source SIEMS out there and a typical SIEM system on average can get a bit complex to understand and so most SIEM's usually are not beginner friendly. The overal goal of this project is design, stimulate a SIEM system just like a real SIEM that can be understood and managed by the common John Doe in which it will allow them to monitor their local machine events. This will also allows us, the developers, as well as the end user to understand event monitoring in the sense of defense and security handling, meaning this will teach us what really goes on the backend of your local machine and how much data truly your system captures and what threats it is constantly trying to prevent and notify as the average Joe is typcially clueless on such matter.
